%!TEX root = ../main.tex
\section{Related Work}
\label{sec:related_work}

\paragraph{Scientific Figure Benchmarks.}
%Recent benchmarks have shifted from synthetic chart QA toward \todo{SE: is it? I thought there were more synthetic recently -- PAUL}real scientific figures and more demanding reasoning tasks. 
\po{Recent benchmarks cover both synthetic and real scientific figures, with growing emphasis on demanding reasoning tasks. ChartQA~\citep{masry2022chartqa} and ChartBench~\citep{xu2024chartbench} evaluate visual and logical reasoning over charts, whereas CharXiv~\citep{wang2024charxiv} and SciFIBench~\citep{huang2024scifibench} focus directly on figures drawn from scientific papers. Other benchmarks broaden this space. ChartMuseum~\citep{tang2025chartmuseum} studies expert-annotated reasoning over real-world charts, ChartQAPro~\citep{masry2025chartqapro} introduces diverse and unanswerable chart questions, EncQA~\citep{mukherjee2025encqa} organises tasks around visual encoding channels, and MultiChartQA~\citep{zhu2024multichartqa} evaluates reasoning across multiple charts.} Table~\ref{tab:related-work-benchmark-comparison} summarises the evaluation scope of existing chart and scientific figure benchmarks. Together, these benchmarks show that scientific figure understanding remains challenging even for frontier VLMs.

\textsc{SciFig-Eval} extends this line of work by jointly evaluating perception, reasoning, and behavioural reliability. Perception is evaluated through open-ended figure descriptions scored using MQM-adapted checklists, reasoning through targeted capability questions, and behaviour through controlled probes that test responses to uncertainty and misleading context. Appendix~\ref{sec:benchmark-comparison} provides a compact comparison of benchmark scope, task type, and evaluation dimensions.


%Recent chart-understanding benchmarks have moved beyond synthetic chart questions toward real charts, scientific figures, and more demanding reasoning tasks. ChartQA and ChartBench evaluate visual and logical reasoning over charts~\citep{masry2022chartqa,xu2024chartbench}, while CharXiv and SciFIBench focus more directly on charts and figures drawn from scientific papers~\citep{wang2024charxiv,huang2024scifibench}. Newer benchmarks further broaden the evaluation space: ChartMuseum tests expert-annotated reasoning over real-world charts, ChartQAPro introduces more diverse and unanswerable chart questions, EncQA organises tasks around visual encoding channels, and MultiChartQA evaluates reasoning across multiple charts~\citep{tang2025chartmuseum,masry2025chartqapro,mukherjee2025encqa,zhu2024multichartqa}. Together, these benchmarks show that chart and figure understanding remains difficult even for frontier VLMs.

\begin{table*}[t]
\centering
\resizebox{\textwidth}{!}{%
\begin{tabular}{lccccccc}
\toprule
\textbf{Benchmark} &
\textbf{Scientific} &
\textbf{Open} &
\textbf{Reasoning} &
\textbf{Robust} &
\textbf{Behaviour} &
\textbf{Uncertainty} &
\textbf{False Premise} \\
\midrule
ChartQA~\citep{masry2022chartqa}
& $\times$ & $\times$ & \checkmark & $\times$ & $\times$ & $\times$ & $\times$ \\

ChartBench~\citep{xu2024chartbench}
& $\times$ & \checkmark & \checkmark & $\times$ & $\times$ & $\times$ & $\times$ \\

MathVista~\citep{lu2024mathvista}
& $\times$ & $\times$ & \checkmark & $\times$ & $\times$ & $\times$ & $\times$ \\

MMMU~\citep{yue2024mmmu}
& Partial & $\times$ & \checkmark & $\times$ & $\times$ & $\times$ & $\times$ \\

CharXiv~\citep{wang2024charxiv}
& \checkmark & \checkmark & \checkmark & $\times$ & $\times$ & $\times$ & $\times$ \\

SciFIBench~\citep{huang2024scifibench}
& \checkmark & \checkmark & \checkmark & Partial & $\times$ & $\times$ & $\times$ \\

ChartMuseum~\citep{tang2025chartmuseum}
& \checkmark & \checkmark & \checkmark & $\times$ & $\times$ & $\times$ & $\times$ \\

ChartQAPro~\citep{masry2025chartqapro}
& $\times$ & $\times$ & \checkmark & \checkmark & $\times$ & Partial & $\times$ \\

\midrule
\textsc{SciFig-Eval} (ours)
& \checkmark & \checkmark & \checkmark & \checkmark & \checkmark & \checkmark & \checkmark \\
\bottomrule
\end{tabular}%
}

\caption{
Comparison of chart and scientific figure benchmarks.
\textsc{SciFig-Eval} uniquely evaluates behavioural reliability under uncertainty and misleading context in addition to perception and reasoning.
}
\label{tab:related-work-benchmark-comparison}
\end{table*} 

%\textsc{SciFig-Eval} builds on this line of work by evaluating scientific figures across three connected dimensions: perception, reasoning, and behaviour. Perception is measured through open-ended figure descriptions scored with MQM-adapted checklists. Reasoning is measured through targeted capability questions. Behaviour is measured through controlled probes that test model responses to uncertainty and misleading context. Appendix~\ref{sec:benchmark-comparison} provides a compact comparison of benchmark scope, task type, and evaluation dimensions.

\paragraph{VLM Hallucination and Reliability.}
Studies show that VLMs can hallucinate or fail under non-standard visual conditions. Prior work on object hallucination and multimodal hallucination benchmarks shows that models can generate fluent yet unsupported visual claims~\citep{rohrbach2018object,li2023pope,guan2024hallusionbench,bai2024hallucination}. In chart settings, CHOCOLATE and ChartHal show that chart captions and answers can contain structured factual errors, including non-existent elements, irrelevant content, and contradictions with the visual evidence~\citep{huang2024chocolate,cui2025charthal}. Robustness benchmarks such as CHAOS \citep{moured2025chaos} and CHART-NOISe \citep{chartdeception2025}
evaluate how chart understanding degrades under perturbation, corruption, and occlusion, while work on misleading visualisations studies how models respond to deceptive chart design~\citep{tonglet2025misleading,chartdeception2025}.

Our admittance dimension relates to calibration research on whether models recognise the limits of their knowledge~\citep{kadavath2022language}. Motivated by this, \textsc{SciFig-Eval} evaluates behaviour under degraded or misleading visual evidence using visual transformations, caption-bias settings, false-premise probes, and selective-blur probes that distinguish unrecoverable missing evidence from contextually inferable evidence. Appendix~\ref{sec:probe-taxonomy} summarises the cognitive motivation of probes.
%Our admittance dimension relates to the broader question of whether models know what they do not know, studied in calibration research where models assign confidence scores to predictions~\citep{kadavath2022language}. These strands motivate evaluating not only whether a model is correct under ordinary conditions, but how it behaves when visual evidence is degraded, incomplete, or contradicted by surrounding text. \textsc{SciFig-Eval} operationalises this setting through visual transformations, figure-in-page conditions, caption-bias probes, false-premise resistance probes, and selective-blur probes that distinguish unrecoverable missing evidence from contextually inferable evidence. Appendix~\ref{sec:probe-taxonomy} summarises the probe families and their cognitive motivations.

\paragraph{Evaluation Methodology.}
The behavioural dimension of \textsc{SciFig-Eval} connects to work on honesty, abstention, sycophancy, and evidence grounding. Benchmarks and surveys on honesty and abstention study whether models know when to answer and when to withhold unsupported claims~\citep{chern2024behonest,wen2024knowlimits,wei2024simpleqa,ren2024selective}. Work on sycophancy and vision-language sycophancy shows that models can align with user-provided or context-provided claims even when those claims conflict with evidence~\citep{sharma2024sycophancy,fanous2025syceval}. \textsc{SciFig-Eval} adapts these concerns to scientific figures by measuring whether models acknowledge uncertainty, resist misleading premises, and draw bounded inferences from partial visual evidence.

\textsc{SciFig-Eval} also builds on structured quality evaluation. We adapt MQM~\citep{lommel2014multidimensional,freitag2021experts} to scientific figure descriptions using checklist-based scoring to evaluate description completeness and error severity.
%\textsc{SciFig-Eval} also builds on structured quality evaluation. MQM decomposes generation quality into typed errors with severity weights~\citep{lommel2014multidimensional,freitag2021experts}; we adapt this approach to scientific figure descriptions and use checklist-based scoring to more accurately evaluate the completeness of figure descriptions. 
To enable the evaluation of model performance at scale, we use LLM judges in the spirit of recent LLM-as-judge work~\citep{zheng2023judging}, but constrain LLM-based judgement to verifiable figure-specific checklist items and validate aggregate rankings against human judgments. 

